\chapter{Contexte général du projet}

\section*{Introduction}
Ce premier chapitre situe le projet dans son contexte. Il présente l'organisme d'accueil, analyse
le processus existant de suivi des projets — fondé sur des fichiers Excel — et en dégage les
limites. De cette étude découlent la problématique, la solution proposée (la plateforme PMS) et la
méthodologie d'ingénierie retenue pour la conduite du projet.

\section{Cadre du projet}
Le présent travail constitue un Projet de Fin d'Études (PFE) en vue de l'obtention du Diplôme
National d'Ingénieur. Il a été réalisé au sein d'une entreprise de services du numérique (ESN), qui
conçoit et met en œuvre des systèmes d'information pour ses clients dans le cadre de projets
contractualisés, le plus souvent au forfait (FP) ou en régie (T\&M).

\section{Organisme d'accueil}
L'entreprise d'accueil opère en tant qu'intégrateur de solutions logicielles. Son activité repose
sur la conduite simultanée de plusieurs projets, mobilisant un \emph{pool} de ressources
(développeurs, analystes métier, chefs de projet) affectées à des projets clients. La rentabilité
de l'entreprise dépend directement de sa capacité à maîtriser, pour chaque projet, l'équilibre
entre la charge vendue, la charge réellement consommée et le chiffre d'affaires facturé.

\section{Étude de l'existant}

\subsection{Le processus actuel}
Le suivi de chaque projet est assuré au moyen d'un classeur Excel — la \og{}~fiche de revue de
projet~\fg{} — comportant une vingtaine d'onglets interconnectés. Ce classeur centralise, pour un
projet donné, l'ensemble des informations de pilotage~: identification et budget du projet, taux de
coût chargé des ressources, plan de charge prévisionnel, charges réelles, jalons de facturation,
missions, risques, livrables et tableaux de bord.

\begin{longtable}{|m{4.5cm}|m{10.5cm}|}
\hline
\textbf{Famille d'onglets} & \textbf{Rôle métier} \\
\hline
\endhead
\endfoot
\endlastfoot
\hline
Identification & Fiche d'identification du projet~: client, contrat, budget, devise, durée,
chef de projet, charge vendue. \\
\hline
Coûts (TCC, Coût Prévisionnel, Coût réel Actualisé) & Calcul du coût journalier des ressources et
des coûts prévisionnels et réels par mois. \\
\hline
Facturation (Avancement, Récapitulatif) & Jalons de facturation en pourcentage du contrat, suivi
des montants facturés et réglés. \\
\hline
Missions & Calcul des coûts de déplacement (perdiem, billet, séjour). \\
\hline
Pilotage (Statistiques, Dashboard, Glossaire) & Consolidation des indicateurs et définitions des
KPI. \\
\hline
Gouvernance (Risques, Livrables, Parties Prenantes, Changements, Avenants) & Registres de suivi de
la vie du projet. \\
\hline
\caption{Structure fonctionnelle du classeur Excel de revue de projet}
\label{tab:excel-structure}
\end{longtable}

\subsection{Critique de l'existant}
Si ce classeur traduit une logique métier riche et éprouvée, son support — le tableur — engendre
des limites structurelles~:

\begin{itemize}
    \item \textbf{Dispersion}~: un classeur par projet, dupliqué et susceptible de diverger d'une
    version à l'autre, sans source unique de vérité.
    \item \textbf{Fiabilité des calculs}~: la dépendance à des formules et à des références entre
    onglets est fragile~; le fichier réel fourni contient d'ailleurs des cellules en erreur
    (\texttt{\#REF!}, \texttt{\#VALUE!}).
    \item \textbf{Visibilité financière}~: l'absence de consolidation rend difficile la lecture du
    coût réel et de la marge à l'échelle du portefeuille de projets.
    \item \textbf{Contrôle d'accès}~: un tableur ne permet pas de cloisonner finement
    l'information — par exemple d'interdire au développeur l'accès aux données financières.
    \item \textbf{Traçabilité}~: l'historique des modifications et des décisions n'est pas conservé
    de manière fiable.
\end{itemize}

\section{Problématique}
La question centrale est donc la suivante~: \emph{comment doter l'entreprise d'un outil centralisé,
fiable et sécurisé qui automatise la logique financière des projets aujourd'hui portée par Excel,
tout en restant suffisamment configurable pour absorber l'évolution des règles métier sans
refonte~?} Deux exigences se dégagent~: préserver la richesse du modèle métier existant (TCC, plan
de charge, EAC, marges, facturation) et garantir une architecture d'autorisation \emph{dynamique},
découplée des rôles.

\section{Solution proposée}
La solution retenue est le développement d'une plateforme web centralisée, \textbf{PMS}
(\textit{Project Management System}). Elle offre une source unique d'information et automatise le
pilotage opérationnel et financier des projets. Ses objectifs sont~:

\begin{itemize}
    \item centraliser les données des projets, ressources et équipes~;
    \item planifier la charge (en jours-homme) et saisir les charges réelles~;
    \item gérer les jalons de facturation et les missions~;
    \item calculer automatiquement les indicateurs financiers (EAC, marges, avancement de
    facturation, reste à faire, dérive)~;
    \item fournir des tableaux de bord décisionnels adaptés à chaque profil d'utilisateur~;
    \item reposer sur un contrôle d'accès dynamique
    \og{}~Rôle~$\rightarrow$~Permission~$\rightarrow$~Module~\fg{}.
\end{itemize}

\section{Méthodologie de travail}
Le projet est conduit selon une démarche d'ingénierie \textbf{par phases successives et validées}~:
analyse métier, conception de l'architecture, conception de la base de données, modélisation UML,
puis réalisation itérative des modules (authentification et RBAC, gestion des utilisateurs, des
rôles et permissions, des projets, des équipes, du TCC, du plan de charge, des charges réelles, de
la facturation, des missions, puis des KPI et du \emph{reporting}). Chaque phase produit des
livrables techniques, met à jour la documentation et alimente le présent rapport.

La conduite du projet s'appuie également sur une \textbf{priorité des sources de vérité}~: les
décisions d'architecture validées priment, suivies des instructions métier, puis de la réalité
extraite du fichier Excel, des spécifications (BRS, cas d'utilisation, SRS) et enfin du plan de
projet, considéré comme une orientation stratégique à challenger.

\section*{Conclusion}
L'étude du contexte et de l'existant a mis en évidence la nécessité de remplacer un suivi Excel
dispersé par une plateforme centralisée, fiable et configurable. Le chapitre suivant approfondit
cette ambition en analysant et en spécifiant précisément les besoins fonctionnels et non
fonctionnels du système, ainsi que les règles de gestion et le modèle financier qui en constituent
le cœur métier.
